\documentclass[11pt,a4paper]{article}

\usepackage[margin=1in]{geometry}
\usepackage{titlesec}
\usepackage{enumitem}
\usepackage{xcolor}
\usepackage{hyperref}
\usepackage{fancyhdr}
\usepackage{booktabs}
\usepackage{array}
\usepackage{multirow}

\definecolor{headerblue}{RGB}{20,60,110}
\definecolor{phasegreen}{RGB}{22,163,74}
\definecolor{phaseorange}{RGB}{234,88,12}
\definecolor{phasepurple}{RGB}{109,40,217}
\definecolor{lightgray}{RGB}{245,245,245}

\titleformat{\section}{\large\bfseries\color{headerblue}}{\thesection}{1em}{}
\titleformat{\subsection}{\normalsize\bfseries\color{phasepurple}}{\thesubsection}{1em}{}
\titleformat{\subsubsection}{\normalsize\bfseries}{\thesubsubsection}{1em}{}

\pagestyle{fancy}
\fancyhf{}
\fancyhead[L]{Kharcha Tracker — R\&D Roadmap v2}
\fancyhead[R]{\thepage}
\renewcommand{\headrulewidth}{0.4pt}

\hypersetup{colorlinks=true, linkcolor=headerblue, urlcolor=headerblue}

\title{%
  \textbf{\Large Kharcha Tracker} \\[4pt]
  \large R\&D Roadmap v2: Completed Phases \& Next Research Directions
}
\author{Umair Abbas \\ AI Engineering Intern}
\date{\today}

% ─────────────────────────────────────────────────────────────────────────────
\begin{document}

\maketitle
\thispagestyle{empty}

% ─────────────────────────────────────────────────────────────────────────────
\section*{Executive Summary}

Kharcha Tracker has progressed through three planned phases of development,
growing from a single-user JSON-based demo into a functional multi-tenant SaaS
prototype. All three phases have been completed and validated locally on a
zero-cost infrastructure stack. This document records what was built, what was
learned during implementation, and proposes new research directions for the next
iteration — including AI-driven insights, automation pipelines, and a path
toward a commercially viable product.

% ─────────────────────────────────────────────────────────────────────────────
\section{Completed Work}

\subsection*{Phase 1 — SaaS Foundation \textnormal{\small\color{phasegreen}(Complete)}}

\begin{itemize}[leftmargin=1.2cm]
  \item Migrated from \texttt{lowdb} (local JSON) to Supabase PostgreSQL with
        row-level security (RLS) for full tenant isolation.
  \item Implemented Supabase JWT authentication — email/password sign-up and
        login, with session persistence via the Supabase client SDK.
  \item Introduced \texttt{workspace\_id} foreign keys on all tables
        (\texttt{expenses}, \texttt{categories}, \texttt{budgets}) so each user's
        data is isolated by workspace.
  \item Defined two roles: \textbf{Owner} (full control) and \textbf{Member}
        (read/write expenses, read budgets). Enforced via RLS policies.
  \item Auto-provisioning: a Supabase trigger creates a default workspace and
        seeds five default categories on every new user sign-up.
  \item Multi-workspace switching in the UI — users can belong to more than one
        workspace and switch context without re-authenticating.
\end{itemize}

\textbf{Key finding:} Supabase RLS is powerful but the interaction between
\texttt{auth.uid()} and service-role clients requires careful separation. All
alert and deduplication logic was moved to the backend with the service-role key
so it could bypass RLS safely without exposing elevated privileges to the
browser.

\subsection*{Phase 2 — Budget Alert Engine \textnormal{\small\color{phasegreen}(Complete)}}

\begin{itemize}[leftmargin=1.2cm]
  \item Event-driven evaluation on every new expense: category-level and
        workspace-level budgets are checked independently.
  \item Three thresholds per budget per month: 80\%, 90\%, 100\%.
        Each combination of \texttt{(workspace\_id, category\_id, month, threshold)}
        fires at most once, enforced by the \texttt{alert\_logs} table with a
        \texttt{UNIQUE NULLS NOT DISTINCT} constraint.
  \item Alert dispatch is fire-and-forget (\texttt{setImmediate}) so it never
        blocks the HTTP response to the expense creation request.
  \item Email delivery via Resend SDK — rich HTML template with workspace name,
        category, spend vs.\ limit, progress bar, and dashboard link.
  \item WhatsApp channel built and feature-flagged
        (\texttt{WHATSAPP\_ENABLED=false}) pending Meta template approval.
  \item A \texttt{public.profiles} table stores optional WhatsApp numbers
        (E.164 format) per user. The \texttt{handle\_new\_user} trigger
        auto-inserts a profile row on sign-up.
  \item Frontend budget banner: shows only the \emph{highest} threshold crossed
        per category (not all three simultaneously). Banners clear immediately
        when a budget is updated.
\end{itemize}

\textbf{Key finding:} Resend's free tier restricts delivery to the account
owner's email unless a custom domain is verified. For multi-user SaaS, domain
verification is a hard requirement before the alert system is commercially
useful.

\subsection*{Phase 3 — AI Receipt Scanning \& UI Polish \textnormal{\small\color{phasegreen}(Complete)}}

\begin{itemize}[leftmargin=1.2cm]
  \item Receipt OCR via Groq API — model migrated from the decommissioned
        \texttt{llama-3.2-90b-vision-preview} to
        \texttt{meta-llama/llama-4-scout-17b-16e-instruct}, which supports
        native JSON mode (\texttt{response\_format: json\_object}) for
        deterministic structured output.
  \item Prompt engineering tuned for Pakistani receipts and mobile wallet
        screenshots (Easypaisa, JazzCash): explicit field mapping for
        ``Total Amount'', Rs.\ prefix, and Pakistani date formats.
  \item Express body limit raised to 10\,MB to accommodate base64-encoded
        receipt images (default 100\,KB limit was silently dropping requests).
  \item Pre-filled form fields highlighted in amber with a
        ``Please verify before saving'' hint — the system never auto-saves.
  \item New abstract geometric SVG logo (\texttt{KharchaLogo}) replaces the
        emoji placeholder; works in both light and dark mode.
  \item Mobile bottom-sheet panel with swipe-to-close, triggered by a floating
        action button (FAB) — visible only on screens below the \texttt{md}
        breakpoint.
  \item TanStack Query (\texttt{@tanstack/react-query}) introduced for the
        expenses list: optimistic insert and optimistic delete so the UI
        responds instantly while the backend confirms in the background.
  \item Budget progress bars animate color transitions:
        green ($<$80\%) $\to$ amber (80--99\%) $\to$ red ($\geq$100\%).
\end{itemize}

\textbf{Key finding:} Groq vision models have high throughput and low latency
but are subject to rapid deprecation cycles (one model was decommissioned during
active development). An abstraction layer around the OCR provider is advisable
so the underlying model can be swapped without touching application logic.

% ─────────────────────────────────────────────────────────────────────────────
\section{Current Technology Stack}

\begin{table}[h]
\centering
\renewcommand{\arraystretch}{1.35}
\begin{tabular}{p{3cm} p{4.5cm} p{5.5cm}}
\toprule
\textbf{Component} & \textbf{Solution} & \textbf{Free-Tier Limit} \\
\midrule
Database \& Auth     & Supabase (PostgreSQL)              & 500\,MB, 50,000 MAU \\
Frontend Framework   & React 18 + Vite + Tailwind CSS     & — \\
State Management     & TanStack Query v5                  & — \\
Charts               & Recharts                           & — \\
Backend              & Node.js + Express (ESM)            & — \\
Frontend Hosting     & Vercel (planned)                   & Unlimited personal projects \\
Backend Hosting      & Render (planned)                   & Free tier, sleeps after 15\,min \\
Email                & Resend                             & 3,000 emails/month \\
WhatsApp             & Meta Cloud API (flagged off)       & 1,000 conversations/month \\
AI OCR               & Groq (Llama 4 Scout)               & Free during preview \\
\bottomrule
\end{tabular}
\caption{Current infrastructure stack}
\end{table}

% ─────────────────────────────────────────────────────────────────────────────
\section{Proposed Research Directions (Phase 4+)}

\subsection{AI-Powered Spending Insights}

\subsubsection*{Monthly Narrative Summary}
Generate a plain-language summary of the user's month: top spending categories,
biggest single expense, comparison to last month, and one actionable
recommendation. Deliverable via email on the 1st of each month (Resend +
cron-job.org) and visible as a collapsible card on the dashboard.

\textit{Research question:} Which LLM prompt structure produces the most
accurate and readable summaries for Pakistani spending contexts (PKR, local
vendor names, utility bills)?

\subsubsection*{Anomaly Detection}
Flag expenses that are statistically unusual relative to the user's own history
(e.g.\ a single Food expense 3$\times$ the monthly average). Can be implemented
as a lightweight z-score check on the backend without a dedicated ML model.

\subsubsection*{Next-Month Spend Prediction}
Use the last 3--6 months of category-level spending to produce a simple linear
forecast. Display as a ``predicted budget'' suggestion when a user opens the
Budget Settings modal for a new month.

\subsection{Automation Pipelines}

\subsubsection*{Bank SMS Parsing}
Pakistani banks (HBL, MCB, UBL, Meezan) send structured SMS alerts on every
transaction. A mobile companion app or browser extension could forward these
SMS messages to \texttt{POST /api/scan-sms}, which parses the text using a
regex + LLM hybrid and auto-creates a draft expense for user confirmation.

\textit{Technical challenge:} Android SMS permissions and iOS restrictions
differ. A PWA cannot read SMS; a native wrapper (Capacitor or React Native)
would be required.

\subsubsection*{Recurring Expense Detection}
Analyse expense history to identify patterns (same vendor, similar amount,
similar date each month) and surface them as suggested entries at the start of
each billing cycle. No user configuration required.

\subsubsection*{Voice Input}
Integrate Groq's Whisper endpoint (\texttt{whisper-large-v3-turbo}) for
voice-to-expense: the user speaks ``Spent 350 on biryani today'' and the app
parses it into a pre-filled form. Groq's Whisper is currently free with no
stated limit.

\subsection{Engagement \& Retention Features}

\subsubsection*{Daily Logging Streaks}
Track consecutive days with at least one expense logged. Display a streak
counter in the header. Streak-break email reminder sent via Resend if no
expense is logged by 9\,PM.

\subsubsection*{Savings Goals}
Users define a goal (e.g.\ ``Save Rs\,20,000 for a laptop by October''). The
system tracks the gap between income and expenses and shows a goal progress bar.
Requires adding an optional income-tracking table.

\subsubsection*{Split Expenses}
Allow a workspace member to flag an expense as shared and assign a percentage
to each member. The system tracks each member's balance and shows who owes whom.
Relevant for roommates and shared household budgets.

\subsection{Data Export \& Reporting}

\subsubsection*{CSV / Excel Export}
A \texttt{GET /api/export?workspace\_id=\&month=\&format=csv} endpoint that
streams a spreadsheet of the filtered expenses. No third-party library needed
for CSV; \texttt{xlsx} (SheetJS) for Excel.

\subsubsection*{Scheduled PDF Report}
Monthly PDF summary (expense table + pie chart + budget status) generated
server-side using \texttt{puppeteer} or \texttt{@react-pdf/renderer} and
delivered via Resend attachment.

\subsection{Infrastructure \& Reliability}

\subsubsection*{Render Cold-Start Mitigation}
Render's free tier suspends the backend after 15\,minutes of inactivity,
causing a 30--50\,second cold start on the next request. Mitigation: register
a free cron at \href{https://cron-job.org}{cron-job.org} to ping
\texttt{/health} every 14\,minutes, keeping the service warm.

\subsubsection*{Provider Abstraction for AI}
Create an \texttt{OcrProvider} interface with a \texttt{scan(dataUri, categories)}
method. Implement two concrete providers: \texttt{GroqProvider} (current) and
\texttt{GeminiProvider} (fallback). The route selects the active provider via an
environment variable (\texttt{OCR\_PROVIDER=groq|gemini}). This insulates the
application from future model deprecations.

\subsubsection*{Audit Log}
Record every expense create, update, and delete in an \texttt{audit\_log} table:
\texttt{(id, workspace\_id, user\_id, action, table\_name, record\_id, old\_data jsonb, new\_data jsonb, created\_at)}.
Essential for team workspaces where multiple members can modify shared records.

\subsection{Monetisation Path}

\begin{table}[h]
\centering
\renewcommand{\arraystretch}{1.35}
\begin{tabular}{p{2.5cm} p{2.5cm} p{7.5cm}}
\toprule
\textbf{Tier} & \textbf{Price} & \textbf{Includes} \\
\midrule
Starter  & Free          & 1 workspace, 5 categories, 3 budget alerts/month, manual entry only \\
Growth   & \$9--15/month & Unlimited categories, unlimited alerts, receipt scanning, CSV export, monthly summary email \\
Team     & \$49+/month   & Multi-member workspaces, audit log, split expenses, scheduled PDF reports, priority support \\
\bottomrule
\end{tabular}
\caption{Proposed monetisation tiers}
\end{table}

% ─────────────────────────────────────────────────────────────────────────────
\section{Suggested Next Sprint (Priority Order)}

\begin{enumerate}[leftmargin=1.2cm]
  \item \textbf{Deploy} — Vercel (frontend) + Render (backend). Config files
        already in repository. Establishes a shareable demo URL for team
        review.
  \item \textbf{CSV Export} — single endpoint, no new dependencies for CSV.
        High user value, low implementation cost.
  \item \textbf{Render keep-alive cron} — 15-minute cron-job.org ping.
        Prevents the most visible user-facing issue on the free tier.
  \item \textbf{Monthly AI Summary} — Groq text model + Resend + cron.
        Differentiates the product from generic expense trackers.
  \item \textbf{OCR Provider Abstraction} — insulates future development from
        Groq model deprecation cycles.
  \item \textbf{Voice Input} — Groq Whisper integration. Fast to build given
        the existing Groq SDK dependency.
\end{enumerate}

% ─────────────────────────────────────────────────────────────────────────────
\section{Conclusion}

The three foundational phases have proven that a production-quality, multi-tenant
expense tracking SaaS can be built entirely on free-tier infrastructure. The
architecture — Supabase for data and auth, Express for backend logic, Groq for
AI features, Resend for notifications — is modular enough that each component
can be upgraded to a paid tier independently as the user base grows. The research
directions proposed in this document focus on features that increase daily
engagement (streaks, voice input, smart summaries) and operational reliability
(provider abstraction, audit logs, cold-start mitigation), both of which are
necessary before a public launch.

\vspace{1cm}
\noindent\rule{\textwidth}{0.4pt}
\small
\noindent Repository: \url{https://github.com/UmairAbbas1/kharcha-tracker} \\
Stack: React 18 · Vite · Tailwind CSS · TanStack Query · Node.js · Express · Supabase · Groq · Resend \\
Document version: 2.0 — \today

\end{document}
